\documentclass[11pt]{article}
\usepackage[a4paper,margin=1in]{geometry}
\usepackage[T1]{fontenc}
\usepackage[utf8]{inputenc}
\usepackage[english]{babel}
\usepackage{lmodern}
\usepackage{microtype}
\usepackage{amsmath,amssymb,mathtools}
\setlength{\parindent}{1.2em}
\setlength{\parskip}{0.35em}
\setlength{\emergencystretch}{5em}
\allowdisplaybreaks
\newcommand{\dd}{\,\mathrm d}
\newcommand{\ii}{\mathrm i}
\newcommand{\cosec}{\operatorname{cosec}}
\newcommand{\tg}{\operatorname{tg}}
\newcommand{\arccosop}{\operatorname{arc\,cos}}
\newcommand{\art}[1]{\bigskip\begin{center}\large\bfseries #1.\end{center}\nopagebreak}
\begin{document}
\begin{center}
{\Large I.\\[0.4em] On the Elements of the Theory of the Eulerian Integrals.}\par
\medskip
{\normalsize Continuation: Articles 9--12.}
\end{center}

\art{9}

A proof essentially different from those so far cited is finally contained in the treatise “Disquisitiones generales circa seriem infinitam etc. Auctore C. F. Gauss.” The whole arrangement of that work, however, makes it impossible to present the proof here, since the foundations on which it rests yield at once a complete theory of the Eulerian integrals, and for that very reason are too substantial to be developed here merely for this single purpose.

To these proofs I now wish to add one which, so far as I know, is new, which requires very little preparation, and which always remains within the domain of integral calculus, even though the method itself is not exactly new: it amounts to reducing the problem to the integration of a differential equation of the second order. In the equation
\[
B B=\int_0^\infty\frac{x^{b-1}\dd x}{x+1}
     \int_0^\infty\frac{y^{b-1}\dd y}{y+1}
=\int_0^\infty\frac{\dd x}{x+1}
     \int_0^\infty\frac{(xy)^{b-1}}{y+1}\dd y
\]
put $y=\frac zx$, $\dd y=\frac{\dd z}{x}$; then reverse the order of integration and carry out the integration with respect to $x$. One then easily finds
\begin{equation}
B B=\int_0^\infty\frac{z^{b-1}\ell z}{z-1}\dd z
=\frac{\dd}{\dd b}\int_0^\infty\frac{z^{b-1}\dd z}{z-1}.
\tag{27}
\end{equation}
By indefinite integration with respect to $b$ one therefore obtains
\[
\int B B\,\dd b=\int_0^\infty\frac{z^{b-1}\dd z}{z-1}.
\]
The integral on the right differs from the integral $B$ only in the form of the denominator, and it is easy to foresee that one recovers the integral $B$ if the expression just obtained is subjected to the same treatment. Indeed, one finds
\[
B\int B B\,\dd b
=\int_0^\infty\frac{\dd z}{z-1}
  \int_0^\infty\frac{(zy)^{b-1}}{y+1}\dd y.
\]
If one now puts $y=\frac xz$, $\dd y=\frac{\dd x}{z}$, then, according to Articles 6 and 7, decomposes the integral taken with respect to $z$ into two integrals whose limits are $0,1-\delta$ and $1+\varepsilon,\infty$, reverses the order of integration, and carries out the integration with respect to $z$, one obtains
\[
B\int B B\,\dd b
=\int_0^\infty\frac{x^{b-1}\ell x}{x+1}\dd x
+\lim\int_0^\infty\frac{x^{b-1}\dd x}{x+1}\,\ell\!\left(\frac{\delta}{\varepsilon}\right).
\]
Here $\lim \ell(\delta/\varepsilon)$ is admittedly undetermined, but in any case it is independent of $x$ and $b$; let it be denoted by $k$. Then the last equation gives
\[
B\int B B\,\dd b=\frac{\dd B}{\dd b}+kB,
\]
or, if one observes that the integral on the left already contains an arbitrary constant,
\begin{equation}
B\int B B\,\dd b=\frac{\dd B}{\dd b}.
\tag{28}
\end{equation}
From this, by division by $B$ and differentiation with respect to $b$, one obtains the second-order differential equation
\[
B B=\frac1B\frac{\dd\dd B}{\dd b^2}
-\frac1{B B}\left(\frac{\dd B}{\dd b}\right)^2.
\]
It has the distinctive property that the independent variable $b$ does not occur; it can therefore, as is well known, be reduced to a differential equation of first order if the first differential quotient is introduced as a new variable. Denote this quotient by $B'$. Then
\[
\frac{\dd B}{\dd b}=B',\qquad
\frac{\dd\dd B}{\dd b^2}=\frac{\dd B'}{\dd b}
=\frac{B'\dd B'}{\dd B}.
\]
If these transformations are introduced into the above differential equation, it becomes
\[
B B=\frac{B'\dd B'}{B\,\dd B}-\frac{B'B'}{B B},
\]
or
\[
\dd(BB)=
\frac{BB\,\dd(B'B')-B'B'\dd(BB)}{(BB)^2}
=\dd\frac{B'B'}{BB}.
\]
Its integral is
\[
BB=cc+\frac{B'B'}{BB}
=cc+\frac1{BB}\left(\frac{\dd B}{\dd b}\right)^2.
\]
To determine the constant, it is now necessary to know two properties of the integral. We take these from Article 3, where it was shown that for $b=\frac12$ the integral $B$ attains a minimum equal to $\pi$; since at the same time $\frac{\dd B}{\dd b}=0$, it follows immediately that $cc=\pi\pi$.

One then obtains further
\[
\pm\,\dd b=
\frac{\dd B}{B\sqrt{BB-\pi\pi}}
=\frac1\pi
  \frac{-\dd\frac{\pi}{B}}
       {\sqrt{1-\frac{\pi\pi}{BB}}},
\]
and therefore, by integration,
\[
\pm(b+c)=\frac1\pi\,\arccosop\frac{\pi}{B},
\qquad
B=\frac{\pi}{\cos(b+c)\pi}.
\]
Since $B$ has the value $\pi$ for $b=\frac12$, one must have
\[
\cos\left(\frac12+c\right)\pi=-\sin c\pi=1,
\qquad \cos c\pi=0.
\]
Consequently
\[
\cos(b+c)\pi=\cos b\pi\cos c\pi-\sin b\pi\sin c\pi=\sin b\pi,
\]
whereby one again finds $B=\frac{\pi}{\sin b\pi}$.

\art{10}

Because the proof given in the preceding article still does not fully satisfy the requirements of the greatest rigor, especially insofar as the integral
\[
\int_0^\infty\frac{x^{b-1}\dd x}{x-1}
\]
plays a role in it, it is not inappropriate to add here a modification in which indefinite integration with respect to $b$ is entirely avoided. Since equation (28) can also be written
\[
\int B B\,\dd b=\frac{\dd\,\ell B}{\dd b},
\]
one may suppose that an expansion of the expression on the right will also lead to the goal. Since $B=\Gamma(b)\Gamma(1-b)$, it is enough to find an integral for $\frac{\dd\,\ell\Gamma(\mu)}{\dd\mu}$, which is known to be obtainable in the following way. From the definition of $\Gamma(\mu)$ it follows that
\[
\frac{\dd\,\ell\Gamma(\mu)}{\dd\mu}
=\frac1{\Gamma(\mu)}
  \int_0^\infty x^{\mu-1}e^{-x}\ell x\,\dd x.
\]
If for $\ell x$ one substitutes the integral
\[
\ell x=\int_0^\infty\frac{e^{-z}-e^{-zx}}{z}\dd z,
\]
which is obtained from the integral $\int_0^1y^{x-1}\dd y=\frac1x$ by integration with respect to $x$ between the limits $1$ and $x$, and by substituting $e^{-z}$ for $y$, then by reversing the order of integration and with the aid of equation (3) one very easily finds
\[
\frac{\dd\,\ell\Gamma(\mu)}{\dd\mu}
=\int_0^\infty\frac{\dd z}{z}
\left(e^{-z}-\frac1{(z+1)^\mu}\right).
\]
Putting $b$ and $(1-b)$ for $\mu$ gives
\[
\frac{\dd\,\ell B}{\dd b}
=\int_0^\infty\frac{\dd z}{z}
\left(\frac1{(z+1)^{1-b}}-\frac1{(z+1)^b}\right),
\]
and if one puts here $z=x-1$ or $z=\frac1x-1$,
\[
\frac{\dd\,\ell B}{\dd b}
=\int_1^\infty\frac{x^b-x^{1-b}}{x-1}\frac{\dd x}{x}
=\int_0^1\frac{x^b-x^{1-b}}{x-1}\frac{\dd x}{x}.
\]
Thus also
\[
2\frac{\dd\,\ell B}{\dd b}
=\int_0^\infty\frac{x^b-x^{1-b}}{x-1}\frac{\dd x}{x}.
\]
Comparison with equation (27) gives at once
\begin{equation}
2\frac{\dd\,\ell B}{\dd b}
=\int_{1-b}^{b} B B\,\dd b
=2\int_{\frac12}^{b} B B\,\dd b,
\tag{29}
\end{equation}
since, by Article 3, $B$ has the same values for $b=\frac12+b'$ and for $b=\frac12-b'$. Differentiating this equation with respect to $b$ gives again the differential equation treated in the preceding article.

\art{11}

Once the proofs just given have put beyond doubt the correctness of the equation
\[
B(r,1-r)=\int_0^\infty\frac{x^{r-1}\dd x}{x+1}
=\frac{\pi}{\sin r\pi},
\]
where $r$ denotes a positive proper fraction, it follows immediately from Article 2 that the Eulerian integrals of the first kind whose two arguments have a positive integer as their sum can actually be represented. One finds
\begin{equation}
B(m+r,n-r)=
\frac{(m+r-1)\cdots(1+r)r\,(n-1-r)\cdots(2-r)(1-r)}
     {(m+n-1)(m+n-2)\cdots2\cdot1}
\frac{\pi}{\sin r\pi}.
\tag{30}
\end{equation}
From this one can also draw an important converse conclusion about the Eulerian integrals of the first kind. Namely, if one puts $m=0$, $n=1$, $r=\frac12$, one finds
\begin{equation}
\Gamma\!\left(\frac12\right)
=\int_0^\infty e^{-x}\frac{\dd x}{\sqrt{x}}
=2\int_0^\infty e^{-xx}\dd x=\sqrt{\pi},
\tag{31}
\end{equation}
and consequently, by equation (4),
\begin{equation}
\Gamma\!\left(n+\frac12\right)
=(2n-1)(2n-3)\cdots5\cdot3\cdot1\,\frac{\sqrt{\pi}}{2^n}.
\tag{32}
\end{equation}

\art{12}

After the cases have been assembled in Articles 2 and 11 in which the Eulerian integrals of both kinds can be represented without the aid of new functions, I shall return once more, in conclusion, to the integral $B$, in order to develop a few further relations between it and other integrals. Among the different forms in which it appears, the following two are quite interesting:
\[
\int_{-\infty}^{+\infty}
\frac{e^{(2b-1)z}+e^{-(2b-1)z}}{e^z+e^{-z}}\dd z
\quad\text{and}\quad
2\int_0^{\frac{\pi}{2}}(\tg\varphi)^{2b-1}\dd\varphi
=2\int_0^{\frac{\pi}{2}}(\tg\varphi)^{1-2b}\dd\varphi .
\]
More important, however, are its generalizations by the introduction of new constants. Among these belongs the integral
\begin{equation}
\int_0^\infty\frac{x^{b-1}\dd x}{x+c}
=\frac{\pi c^{b-1}}{\sin b\pi},
\tag{33}
\end{equation}
where $c$ must be positive; yet it can subsequently be proved that $c$ may also be imaginary. Namely, if one multiplies numerator and denominator of the function $\frac{x^{b-1}}{x+i}$ by $(x-i)$, then the corresponding integral splits into two others, which can be reduced to the integral $B$ by the substitution $xx=y$. Thus one finds the validity of equation (33) for imaginary $c$. By differentiation and integration with respect to $c$, one can then again derive a series of other integrals.

If $c_1,c_2,\ldots,c_n$ are mutually distinct positive or imaginary constants, further if $0<b<n$, and if
\[
f(x)=(x+c_1)(x+c_2)\cdots(x+c_n),\qquad
f'(x)=\frac{\dd f(x)}{\dd x},
\]
then one also finds
\begin{equation}
\int_0^\infty\frac{x^{b-1}\dd x}{f(x)}
=\frac{\pi}{\sin b\pi}
 \sum_{k=1}^{n}\frac{c_k^{\,b-1}}{f'(-c_k)},
\tag{34}
\end{equation}
by decomposing $\frac{x^\beta}{f(x)}$ into partial fractions, where $\beta$ denotes the greatest integer contained in $b$.

A further advantage of the proof given in Article 10 is that it directly teaches how to determine a related class of integrals. From equations (27) and (29) it follows immediately that
\[
\int_0^\infty\frac{x^b-x^{\frac12}}{1-x}\frac{\dd x}{x}
=\pi\cot b\pi,
\]
and consequently also
\begin{equation}
\int_0^\infty\frac{x^a-x^b}{1-x}\frac{\dd x}{x}
=\pi(\cot a\pi-\cot b\pi).
\tag{35}
\end{equation}
From this one also obtains the integral denoted by $\varphi(b)$ in Article 3:
\[
\begin{aligned}
\int_0^\infty\frac{1-x^\mu}{1-x^\nu}x^{b-1}\dd x
&=\frac1\nu\int_0^\infty
  \frac{x^{\frac b\nu}-x^{\frac{\mu+b}{\nu}}}{1-x}\frac{\dd x}{x}\\
&=\frac{\pi}{\nu}
  \frac{\sin\frac{\mu\pi}{\nu}}
       {\sin\frac{b\pi}{\nu}\sin\frac{(\mu+b)\pi}{\nu}}\\
&=\frac{2\pi}{\nu}
  \frac{\sin\frac{\mu\pi}{\nu}}
       {\cos\frac{\mu\pi}{\nu}-\cos\frac{(\mu+2b)\pi}{\nu}},
\end{aligned}
\]
and its minimum:
\[
\int_0^\infty
\frac{x^{\frac{\mu}{2}}-x^{-\frac{\mu}{2}}}
     {x^{\frac{\nu}{2}}-x^{-\frac{\nu}{2}}}
\frac{\dd x}{x}
=\frac{2\pi}{\nu}\,\tg\frac{\mu\pi}{2\nu}.
\]

Most of the integrals belonging here are found in the “Sammlung von Integraltafeln” issued by Minding at the commission of the Prussian ministry (Berlin, 1849), at the beginning of the fourth division. Perhaps this is the place to indicate and correct several errors found there. By decomposing $\frac1{(x-1)(x+c)}$ into partial fractions and applying the formulas of this article, one easily finds
\[
\int_0^\infty\frac{x^a-1}{(x-1)(x+c)}\dd x
=\frac{\pi}{c+1}
\left(\frac{c^a-\cos a\pi}{\sin a\pi}
-\frac{\ell c}{\pi}\right).
\]
This equation holds for positive and negative proper fractional $a$, as one easily verifies by writing $\frac1x$ and $\frac1c$ instead of $x$ and $c$. If one differentiates with respect to $a$ and then puts $a=0$, one obtains a series of integrals that are given incorrectly in those tables. To make the tedious differentiations as easy as possible, the following device may be used. Put
\[
\frac{\ell c}{\pi}+\frac{c+1}{\pi}
\int_0^\infty\frac{x^a-1}{(x-1)(x+c)}\dd x
=\frac{c^a-\cos a\pi}{\sin a\pi}=J,
\qquad
\frac{\dd^nJ}{\dd a^n}={}^{n}J.
\]
Then for $a=0$ one obtains
\[
\int_0^\infty\frac{(\ell x)^n\dd x}{(x-1)(x+c)}
=\frac{\pi}{c+1}\,{}^{n}J_0.
\]
But one easily finds
\[
\begin{aligned}
&J\pi^n\sin\left(a+\frac n2\right)\pi
+n\,{}^{1}J\,\pi^{n-1}\sin\left(a+\frac{n-1}{2}\right)\pi+\cdots\\
&\quad{}+n\,{}^{n-1}J\,\pi\sin\left(a+\frac12\right)\pi
+{}^{n}J\sin a\pi\\
&=\frac{\dd^n(J\sin a\pi)}{\dd a^n}
=c^a(\ell c)^n-\pi^n\cos\left(a+\frac n2\right)\pi.
\end{aligned}
\]
For $a=0$ this gives recursion formulas for the $J_0$ with even indices and those with odd indices. Thus one finds
\[
\begin{aligned}
\int_0^\infty\frac{\ell x\,\dd x}{(x-1)(x+c)}
&=\frac{(\ell c)^2+\pi^2}{2(c+1)},\\
\int_0^\infty\frac{(\ell x)^2\dd x}{(x-1)(x+c)}
&=\frac{\ell c\bigl((\ell c)^2+\pi^2\bigr)}{3(c+1)},\\
\int_0^\infty\frac{(\ell x)^3\dd x}{(x-1)(x+c)}
&=\frac{\bigl((\ell c)^2+\pi^2\bigr)^2}{4(c+1)},\\
\int_0^\infty\frac{(\ell x)^4\dd x}{(x-1)(x+c)}
&=\frac{\ell c\bigl((\ell c)^2+\pi^2\bigr)
        \bigl(3(\ell c)^2+7\pi^2\bigr)}{5(c+1)\cdot3},\\
\int_0^\infty\frac{(\ell x)^5\dd x}{(x-1)(x+c)}
&=\frac{\bigl((\ell c)^2+\pi^2\bigr)^2
        \bigl((\ell c)^2+3\pi^2\bigr)}{6(c+1)}.
\end{aligned}
\]
In those tables, however, in place of the divisors $2,3,4,5,6$, the divisors $1\cdot2$, $1\cdot2\cdot3$, $1\cdot2\cdot3\cdot4$, $1\cdot2\cdot3\cdot4\cdot5$, $1\cdot2\cdot3\cdot4\cdot5\cdot6$ are given.

\end{document}
